\chapter{Complexity — Uniform Tableau Spec}
%%% generated by the blueprint pipeline from TCSlib.Complexity.CircuitComplexity.UniformTableauSpec
\section{Overview}

Declarations of the Lean module \texttt{TCSlib.Complexity.CircuitComplexity.UniformTableauSpec}.

\section{Declarations}

\begin{definition}[Input-bit layout]
\lean{Complexity.tabLayout}
\label{Complexity.tabLayout}
\leanok
\uses{BoolCircuit.BitSrc}
For a nonnegative integer $N$, define the input-bit layout to be the list of input-bit
references $0,1,\ldots,N-1$, in that order.
\end{definition}

\begin{theorem}[Length of the input layout]
\lean{Complexity.length_tabLayout}
\label{Complexity.length_tabLayout}
\leanok
\uses{Complexity.tabLayout}
\difficulty{1}
For every natural number $N$, the list of input-bit sources in the layout for an input
of length $N$ has length $N$.
\end{theorem}

\begin{theorem}[Validity of the tableau input layout]
\lean{Complexity.tabLayout_valid}
\label{Complexity.tabLayout_valid}
\leanok
\uses{BoolCircuit.BitSrc.Valid, Complexity.tabLayout}
\difficulty{4}
For all natural numbers $N,w$, every bit source in the tableau layout of length $N$ is
valid for an input of length $N$, block width $w$, and instruction index $0$.
\end{theorem}

\begin{theorem}[Evaluation of the tab layout]
\lean{Complexity.tabLayout_eval}
\label{Complexity.tabLayout_eval}
\leanok
\uses{BoolCircuit.BitSrc.eval, Complexity.tabLayout}
\difficulty{3}
For every finite Boolean list $v$, evaluating the tab layout of length $|v|$ on input
$v$, with no previously computed blocks, yields exactly $v$.
\end{theorem}

\begin{definition}[Tabulation configuration circuit]
\lean{Complexity.cfgTab}
\label{Complexity.cfgTab}
\leanok
\uses{BoolCircuit.DAGCircuit, BoolCircuit.tabCircuit, Complexity.cfgArity, Complexity.cfgF, Complexity.cfgProg, Complexity.cfgWidth, Complexity.snapIdx, Complexity.snapWidth, Complexity.tabLayout}
Let $M$ be a finite-tape Turing machine, let $n,T$ be nonnegative integers, and let $z$
be a finite Boolean word. Define the tabulation configuration circuit on $n$ input bits
by applying the tabulation-circuit construction with the parameters derived from $M$,
with input width $|z|+n$, the input-bit layout of that width, step budget $T$, and the
snapshot index and width specified by the machine-dependent snapshot functions at $T$.
\end{definition}

\begin{theorem}[Configuration tableau evaluation]
\lean{Complexity.cfgTab_eval}
\label{Complexity.cfgTab_eval}
\leanok
\uses{BoolCircuit.BitSrc.eval, BoolCircuit.DAGCircuit.eval, BoolCircuit.tabCircuit_eval, Complexity.cfgExp_snapIdx, Complexity.cfgProg, Complexity.cfgProg_blocks, Complexity.cfgProg_wf, Complexity.cfgTab, Complexity.cfgWidth, Complexity.emitted, Complexity.length_cfgProg, Complexity.length_tabLayout, Complexity.snapExp, Complexity.snapIdx, Complexity.snapIdx_lt_succ, Complexity.snapshotAt, Complexity.tabLayout, Complexity.tabLayout_eval, Complexity.tabLayout_valid, Complexity.tableauAccepted}
\difficulty{5}
Let $M$ be a finite-state Turing machine, and let $n,T$ be natural numbers, $z$ a
Boolean list, and $x:\{0,1\}^n\to\{0,1\}$ an input assignment. The configuration tableau
circuit on input $x$ outputs true exactly when, at some time $s<T$, the transition from
the configuration of $M$ on input $z$ followed by the $n$ bits of $x$ emits true.
\end{theorem}

\begin{theorem}[Configuration tableau fan-in bound]
\lean{Complexity.cfgTab_isFaninTwo}
\label{Complexity.cfgTab_isFaninTwo}
\leanok
\uses{BoolCircuit.DAGCircuit.IsFaninTwo, BoolCircuit.tabCircuit_isFaninTwo, Complexity.cfgTab}
\difficulty{1}
Let $M$ be a finite-state Turing machine, and let $n,T$ be natural numbers and $z$ a
Boolean list. The configuration tableau circuit is well formed, and every gate in it
reads at most two vertices.
\end{theorem}

\begin{theorem}[Size of the configuration tableau circuit]
\lean{Complexity.cfgTab_size}
\label{Complexity.cfgTab_size}
\leanok
\uses{BoolCircuit.DAGCircuit.size, BoolCircuit.tabCircuit_size, BoolCircuit.tabW, Complexity.cfgArity, Complexity.cfgF, Complexity.cfgTab, Complexity.cfgWidth, Complexity.length_cfgProg, Complexity.length_tabLayout}
\difficulty{3}
Let $M$ be a finite deterministic Turing machine over the binary alphabet, let
$n,T\in\mathbb{N}$, and let $z$ be a finite binary word. The size of the configuration
tableau circuit with $n$ input variables, initial word $z$, and time bound $T$ is
$n+2+|z|+W\bigl((T+1)(k(2T+1)+(|z|+n)+3)+1\bigr)$, where $k$ is the number of work tapes
of $M$ and $W$ is the fixed instruction-block width of the tableau construction for $M$.
\end{theorem}

\begin{definition}[First residue index]
\lean{Complexity.UTab.rNC}
\label{Complexity.UTab.rNC}
\leanok
The distinguished residue index $r_{\rm NC}$ is the element $0$ of $	ext{Fin}(11)$, the
set of integers from $0$ through $10$.
\end{definition}

\begin{definition}[Second residue index]
\lean{Complexity.UTab.rZ}
\label{Complexity.UTab.rZ}
\leanok
The distinguished residue index $r_{\rm Z}$ is the element $1$ of $	ext{Fin}(11)$, the
set of integers from $0$ through $10$.
\end{definition}

\begin{definition}[Third residue index]
\lean{Complexity.UTab.rT}
\label{Complexity.UTab.rT}
\leanok
The distinguished residue index $r_{\rm T}$ is the element $2$ of $	ext{Fin}(11)$, the
set of integers from $0$ through $10$.
\end{definition}

\begin{definition}[Fourth residue index]
\lean{Complexity.UTab.rI}
\label{Complexity.UTab.rI}
\leanok
The distinguished residue index $r_{\rm I}$ is the element $3$ of $	ext{Fin}(11)$, the
set of integers from $0$ through $10$.
\end{definition}

\begin{definition}[Fifth residue index]
\lean{Complexity.UTab.rQ}
\label{Complexity.UTab.rQ}
\leanok
The distinguished residue index $r_{\rm Q}$ is the element $4$ of $	ext{Fin}(11)$, the
set of integers from $0$ through $10$.
\end{definition}

\begin{definition}[Sixth residue index]
\lean{Complexity.UTab.rLB}
\label{Complexity.UTab.rLB}
\leanok
The distinguished residue index $r_{\rm LB}$ is the element $5$ of $	ext{Fin}(11)$, the
set of integers from $0$ through $10$.
\end{definition}

\begin{definition}[Seventh residue index]
\lean{Complexity.UTab.rLT}
\label{Complexity.UTab.rLT}
\leanok
The distinguished residue index $r_{\rm LT}$ is the element $6$ of $	ext{Fin}(11)$, the
set of integers from $0$ through $10$.
\end{definition}

\begin{definition}[Eighth residue index]
\lean{Complexity.UTab.rLL}
\label{Complexity.UTab.rLL}
\leanok
The distinguished residue index $r_{\rm LL}$ is the element $7$ of $	ext{Fin}(11)$, the
set of integers from $0$ through $10$.
\end{definition}

\begin{definition}[Register index eight]
\lean{Complexity.UTab.rPA}
\label{Complexity.UTab.rPA}
\leanok
For any finite-state Turing machine $M$ over the binary alphabet, the register index
denoted by $r_{PA}$ is $8$ in the set $	ext{Fin }11$.
\end{definition}

\begin{definition}[Register index nine]
\lean{Complexity.UTab.rPB}
\label{Complexity.UTab.rPB}
\leanok
For any finite-state Turing machine $M$ over the binary alphabet, the register index
denoted by $r_{PB}$ is $9$ in the set $	ext{Fin }11$.
\end{definition}

\begin{definition}[Register index ten]
\lean{Complexity.UTab.rTMP}
\label{Complexity.UTab.rTMP}
\leanok
For any finite-state Turing machine $M$ over the binary alphabet, the register index
denoted by $r_{TMP}$ is $10$ in the set $	ext{Fin }11$.
\end{definition}

\begin{definition}[Sources for tableau bits]
\lean{Complexity.UTab.SSrc}
\label{Complexity.UTab.SSrc}
\leanok
For each natural number $k$, define a source type whose elements are constants carrying
a Boolean value; the input sources $z$ and $u$; previous-configuration sources indexed
by $d\in\{0,1,2\}$ and a bit position $j\in\mathbb N$; current-configuration and
position sources indexed by $j\in\mathbb N$; or end-of-tape sources indexed by
$\tau\in\{0,\ldots,k-1\}$ and $j\in\mathbb N$.
\end{definition}

\begin{definition}[Conversion to circuit bit sources]
\lean{Complexity.UTab.SSrc.toBitSrc}
\label{Complexity.UTab.SSrc.toBitSrc}
\leanok
\uses{BoolCircuit.BitSrc, Complexity.UTab.SSrc, Complexity.UTab.rI, Complexity.UTab.rLB, Complexity.UTab.rPA, Complexity.UTab.rPB, Complexity.UTab.rQ, Complexity.UTab.rT, Complexity.UTab.rZ}
Let $k\in\mathbb N$, let $\rho$ assign a natural number to each of the eleven registers,
and let $s$ be a tableau-bit source. Its corresponding circuit bit source is given as
follows: a constant remains that Boolean constant; $z$ becomes input bit $\rho(8)$; $u$
becomes input bit $\rho(1)+\rho(9)$; a previous-configuration bit $(d,j)$ becomes block
bit $j$ of block $\rho(4)+d-1$; a current-configuration bit $j$ becomes block bit $j$ of
block $\rho(3)-1$; a position bit $j$ becomes block bit $j$ of block $\rho(5)-1$; and an
end-of-tape bit $(\tau,j)$ becomes block bit $j$ of block
$\rho(5)+(\tau+1)(2\rho(2)+1)-1$. All arithmetic in block indices is natural-number
arithmetic, with subtraction truncated at zero.
\end{definition}

\begin{definition}[Validity conditions for tableau sources]
\lean{Complexity.UTab.SSrc.Ok}
\label{Complexity.UTab.SSrc.Ok}
\leanok
\uses{Complexity.UTab.SSrc, Complexity.UTab.rI, Complexity.UTab.rLB, Complexity.UTab.rPA, Complexity.UTab.rQ, Complexity.UTab.rZ}
Let $k\in\mathbb N$, let $\rho$ assign a natural number to each of the eleven registers,
and let $s$ be a tableau-bit source. The source is valid if it is $z$ and
$\rho(8)<\rho(1)$, or if it is a previous-configuration source indexed by $d$ and $j$
and $1\leq\rho(4)+d$. It is also valid if it is a current-configuration source and
$1\leq\rho(3)$, or a position source and $1\leq\rho(5)$; constants, $u$, and end-of-tape
sources are always valid.
\end{definition}

\begin{definition}[Linear expressions for tableau sources]
\lean{Complexity.UTab.SSrc.toLinE}
\label{Complexity.UTab.SSrc.toLinE}
\leanok
\uses{BoolCircuit.tabG, BoolCircuit.tabW, Complexity.CounterProg.LinE, Complexity.UTab.SSrc, Complexity.UTab.rI, Complexity.UTab.rLB, Complexity.UTab.rNC, Complexity.UTab.rPA, Complexity.UTab.rPB, Complexity.UTab.rQ, Complexity.UTab.rT, Complexity.UTab.rZ, Complexity.cfgArity, Complexity.cfgF, Complexity.cfgWidth, Turing.FinTM}
Let $M$ be a finite-state Turing machine over the binary alphabet with $k$ work tapes.
Write $A$ for its configuration arity, $w$ for its configuration width, $G$ for the
maximum width of its transition-table gadgets, and $W=A+G+w$. Each tableau source is
assigned a linear expression over eleven registers, given as a list of register indices
followed by a natural-number constant. A constant source $b$ is assigned $([0],1)$ if
$b$ is true and $([0],0)$ otherwise; the zero-input source is $([0,8],2)$, and the
unit-input source is $([9],0)$. A previous-configuration source $(d,j)$ is assigned the
list consisting of $0$, then $1$, then $W$ copies of $4$, with constant $2+Wd+A+G+j$; a
current-configuration source $j$ and a position source $j$ are each assigned the list
consisting of $0$, then $1$, then $W$ copies of $3$, with constant $2+A+G+j$. An
end-of-tape source $(\tau,j)$ is assigned the list consisting of $0$, then $1$, then $W$
copies of $5$, followed by $2W(\tau+1)$ copies of $2$, with constant
$2+W(\tau+1)+A+G+j$.
\end{definition}

\begin{theorem}[Tabulated source index]
\lean{Complexity.UTab.tabSrc_toBitSrc}
\label{Complexity.UTab.tabSrc_toBitSrc}
\leanok
\uses{BoolCircuit.BitSrc.Valid, BoolCircuit.tabBase, BoolCircuit.tabSrc, Complexity.CounterProg.LinE.val, Complexity.UTab.SSrc, Complexity.UTab.SSrc.Ok, Complexity.UTab.SSrc.toBitSrc, Complexity.UTab.SSrc.toLinE, Complexity.UTab.rI, Complexity.UTab.rLB, Complexity.UTab.rNC, Complexity.UTab.rQ, Complexity.UTab.rT, Complexity.UTab.rZ, Complexity.cfgArity, Complexity.cfgF, Complexity.cfgWidth}
\difficulty{5}
Let $M$ be a finite-state Turing machine with Boolean symbols and $k$ work tapes, let
$n,i\in\mathbb{N}$, let $z$ be a finite bit string, and let $\rho$ assign a natural
number to each of eleven registers. Suppose $\rho(r_{\mathrm{NC}})=n$ and
$\rho(r_{\mathrm{Z}})=|z|$. Let $s$ be a source descriptor, and write $b_s$ for the bit
source obtained from $s$ and $\rho$, and $e_s$ for the linear register expression
associated with $s$ and $M$. Assume that $s$ is available: if it is an input-head
source, then $\rho(r_{\mathrm{PA}})<\rho(r_{\mathrm{Z}})$; if it is a previous-block
source with direction $d\in\{0,1,2\}$, then $1\leq\rho(r_{\mathrm{Q}})+d$; if it is a
current-block source, then $1\leq\rho(r_{\mathrm{I}})$; and if it is a previous-snapshot
source, then $1\leq\rho(r_{\mathrm{LB}})$. Constant, unary-input, and cell-end sources
are always available. If $b_s$ is valid for input length $|z|+n$, output width
$\operatorname{cfgWidth}(M)$, and instruction index $i$, then the tabulated index of
$b_s$ equals the value of $e_s$ under $\rho$.
\end{theorem}

\begin{definition}[Conditional previous-configuration source]
\lean{Complexity.UTab.nbS}
\label{Complexity.UTab.nbS}
\leanok
\uses{Complexity.UTab.SSrc}
For a natural number $k$, a Boolean value $b$, an index $d \in \mathrm{Fin}(3)$, and $j
\in \mathbb{N}$, define $\mathrm{nbS}(b,d,j)$ to be the previous-configuration source
indexed by $(d,j)$ if $b$ is true, and the constant-false source otherwise.
\end{definition}

\begin{definition}[Conditional current source]
\lean{Complexity.UTab.chS}
\label{Complexity.UTab.chS}
\leanok
\uses{Complexity.UTab.SSrc}
For a natural number $j$ and a Boolean $b$, define a source to be the current source at
position $j$ when $b$ is true, and the constant false source when $b$ is false.
\end{definition}

\begin{definition}[Conditional postponed sources]
\lean{Complexity.UTab.psnS}
\label{Complexity.UTab.psnS}
\leanok
\uses{Complexity.UTab.SSrc}
For a natural number $W$ and a Boolean $b$, form a list of length $W$, indexed by
$j=0,\ldots,W-1$. Its $j$th entry is the constant false source if $b$ is true, and the
postponed source at position $j$ if $b$ is false.
\end{definition}

\begin{definition}[Padding a source list]
\lean{Complexity.UTab.padS}
\label{Complexity.UTab.padS}
\leanok
\uses{Complexity.UTab.SSrc}
For a natural number $m$ and a list $l$ of sources, append $\max(m-|l|,0)$ copies of the
constant false source to $l$, where $|l|$ is the length of the list.
\end{definition}

\begin{definition}[Cell source list]
\lean{Complexity.UTab.cellS}
\label{Complexity.UTab.cellS}
\leanok
\uses{Complexity.UTab.SSrc, Complexity.UTab.chS, Complexity.UTab.nbS, Complexity.UTab.psnS}
For a natural number $W$ and Booleans $t_0,r_{\mathrm{lo}},r_T,r_{\mathrm{hi}}$, define
a source list beginning with the constant sources $t_0$ and $r_T$, followed by the nine
neighbor-source terms specified by the conditions $\neg t_0\land r_{\mathrm{lo}}$, $\neg
t_0$, and $\neg t_0\land r_{\mathrm{hi}}$ at the indicated indices. Append the two
conditional current sources at positions $3$ and $4$, each selected using
$r_{\mathrm{lo}}$, and then append the conditional postponed-source list of length $W$
determined by $t_0$.
\end{definition}

\begin{definition}[Layout source]
\lean{Complexity.UTab.layS}
\label{Complexity.UTab.layS}
\leanok
\uses{Complexity.UTab.SSrc}
For each position $i\in\{0,1,2\}$, assign the constant false source when $i=0$, the
input source when $i=1$, and the auxiliary input source when $i=2$.
\end{definition}

\begin{definition}[Input source list]
\lean{Complexity.UTab.inpS}
\label{Complexity.UTab.inpS}
\leanok
\uses{Complexity.UTab.SSrc, Complexity.UTab.chS, Complexity.UTab.layS, Complexity.UTab.nbS, Complexity.UTab.psnS}
For a natural number $W$, Booleans $t_0,p_0,p_1,p_{N1}$, and a layout position
$i\in\{0,1,2\}$, form a list beginning with constant sources carrying, in order,
$t_0,p_1,p_0,p_{N1}$, and $\neg p_0\land\neg p_{N1}$, followed by the layout source at
$i$. Add the three neighbor-source terms specified by $\neg t_0\land\neg p_0$, $\neg
t_0$, and $\neg t_0\land\neg p_{N1}$ at their indicated indices, then the conditional
current sources at positions $1$ and $2$, each selected using $\neg p_0$, and finally
the conditional postponed-source list of length $W$ determined by $t_0$.
\end{definition}

\begin{definition}[Snapshot source list]
\lean{Complexity.UTab.snapS}
\label{Complexity.UTab.snapS}
\leanok
\uses{Complexity.UTab.SSrc, Complexity.UTab.psnS}
For a natural number $W$, a Boolean $t_0$, and a natural number $k$, form a list
beginning with the constant source carrying $t_0$, followed by the constant false source
if $t_0$ is true and otherwise the postponed source at position $W$, and then the
current sources at positions $1$ and $2$. For each $\tau\in\{0,\ldots,k-1\}$, append the
two end sources indexed by $\tau$ at positions $3$ and $4$; finally append the
conditional postponed-source list of length $W$ determined by $t_0$.
\end{definition}

\begin{theorem}[Padding commutes with source translation]
\lean{Complexity.UTab.padS_toBitSrc}
\label{Complexity.UTab.padS_toBitSrc}
\leanok
\uses{Complexity.UTab.SSrc, Complexity.UTab.SSrc.toBitSrc, Complexity.UTab.padS}
\difficulty{1}
Let $k,m$ be natural numbers, let $l$ be a list of sources for $k$ tapes, and let $\rho$
assign natural numbers to the eleven registers. Translate each source to its associated
Boolean circuit bit source. Translating the list obtained by padding $l$ to length at
least $m$ with false constants gives the translated list $l$, followed by $m-|l|$ false
constants.
\end{theorem}

\begin{theorem}[Previous-snapshot source translation]
\lean{Complexity.UTab.psnS_toBitSrc}
\label{Complexity.UTab.psnS_toBitSrc}
\leanok
\uses{BoolCircuit.BitSrc, Complexity.UTab.SSrc.toBitSrc, Complexity.UTab.psnS, Complexity.UTab.rLB, Complexity.cfgStride, Complexity.prevSnapSrcs, Complexity.snapIdx, Complexity.snapWidth}
\difficulty{4}
Let $M$ be a finite Turing machine over the Boolean alphabet, and let $N,T,t$ be natural
numbers. Let $\ell$ be a list of bit sources of length $N$, and assign a natural number
to each of the eleven registers, with register $LB$ equal to $t$ times the configuration
stride for $M$, input length $N$, and time bound $T$. Translating the snapshot source
list of width equal to the number of possible snapshots of $M$, using the flag that is
true exactly when $t=0$, yields the list of sources for the preceding snapshot at time
$t$.
\end{theorem}

\begin{theorem}[Cell source translation]
\lean{Complexity.UTab.cellS_toBitSrc}
\label{Complexity.UTab.cellS_toBitSrc}
\leanok
\uses{BoolCircuit.BitSrc, Complexity.UTab.SSrc.toBitSrc, Complexity.UTab.cellS, Complexity.UTab.chS, Complexity.UTab.nbS, Complexity.UTab.psnS_toBitSrc, Complexity.UTab.rI, Complexity.UTab.rLB, Complexity.UTab.rQ, Complexity.cellIdx, Complexity.cellSrcs, Complexity.cfgStride, Complexity.chSrc, Complexity.nbSrc, Complexity.snapWidth}
\difficulty{5}
Let $M$ be a finite Turing machine over the binary alphabet with $k$ work tapes, and let
$N,T,t,r$ be natural numbers, $\tau$ a work-tape index, and $\rho$ an assignment of
natural numbers to the eleven registers. Let $\sigma$ be a list of bit sources of length $N$.
Suppose the instruction register gives the index of cell $(t,\tau,r)$, the preceding-cell
register gives the index of cell $(t-1,\tau,r)$ whenever $t\ne 0$, and the
first-instruction register equals $t$ times the configuration stride. Then translating
the cell source list for time $t$ and position $r$ through $\rho$ yields exactly the
bit-source list for cell $(t,\tau,r)$, with the specified constants and neighbor,
current-cell, and snapshot sources in their prescribed order.
\end{theorem}

\begin{theorem}[Translation of input-position sources]
\lean{Complexity.UTab.inpS_toBitSrc}
\label{Complexity.UTab.inpS_toBitSrc}
\leanok
\uses{BoolCircuit.BitSrc, Complexity.UTab.SSrc.toBitSrc, Complexity.UTab.chS, Complexity.UTab.inpS, Complexity.UTab.layS, Complexity.UTab.nbS, Complexity.UTab.psnS_toBitSrc, Complexity.UTab.rI, Complexity.UTab.rLB, Complexity.UTab.rQ, Complexity.cfgStride, Complexity.ichSrc, Complexity.ihSrc, Complexity.inpIdx, Complexity.inpSrcs, Complexity.snapWidth}
\difficulty{5}
Let $M$ be a finite Turing machine with $k$ work tapes, and let $N,T,t,p\in\mathbb N$.
Let $\ell$ be a list of bit-source descriptors with $|\ell|=N$, and assume $p<N+2$. For
a valuation $\rho$ of the eleven registers, assume its $I$ register equals the encoded
index of input position $p$ at time $t$, its $Q$ register equals the encoded index of
position $p$ at time $t-1$ whenever $t\ne0$, and its $LB$ register equals $t$ times the
configuration stride. Choose a layout value among $0,1,2$, and assume its associated
source translates to the source from $\ell$ at position $p$ when $1\le p\le N$, and to
the false constant otherwise. Then translating each descriptor in the input-position
source list (formed using the snapshot width, $t$, $p$, and this layout value) gives
exactly the canonical source list for input position $p$ at time $t$, including its
neighboring input and previous-snapshot sources.
\end{theorem}

\begin{theorem}[Snapshot source translation]
\lean{Complexity.UTab.snapS_toBitSrc}
\label{Complexity.UTab.snapS_toBitSrc}
\leanok
\uses{BoolCircuit.BitSrc, Complexity.UTab.SSrc.toBitSrc, Complexity.UTab.psnS_toBitSrc, Complexity.UTab.rI, Complexity.UTab.rLB, Complexity.UTab.rT, Complexity.UTab.snapS, Complexity.cellIdx, Complexity.cfgStride, Complexity.inpIdx, Complexity.snapIdx, Complexity.snapSrcs, Complexity.snapWidth}
\difficulty{4}
Let $M$ be a finite Turing machine over the binary alphabet with $k$ work tapes, and let
$N,T,t\in\mathbb{N}$. Let $\ell$ be a list of Boolean circuit bit sources of length $N$,
and let $\rho$ assign a natural number to each of the 11 registers. Suppose the register
for the current snapshot index has value $\operatorname{snapIdx}_M(N,T,t)$, the register
for the time bound has value $T$, and the register for the left boundary has value
$t\,\operatorname{cfgStride}_M(N,T)$. Then translating each entry of the snapshot-source
list of width equal to the number of possible snapshots, with its initial-time flag set
to the indicator of $t=0$, gives exactly the list of bit sources specified by the
tableau layout at time $t$.
\end{theorem}
